\documentclass[12pt]{article}
% Préambule commun à tous les documents extraits de « La Revue CAPES »
\usepackage[T1]{fontenc}
\usepackage[utf8]{inputenc}
\usepackage{lmodern}
\usepackage{amsmath,amssymb,amsthm,amsfonts,mathrsfs}
\usepackage{setspace}
\usepackage{listings}
\usepackage{pgfplots}
\pgfplotsset{compat=1.18}
\usepackage{longtable}
\usepackage{adjustbox}
\usepackage{lettrine}
\usepackage{tkz-tab}
\usepackage{changepage}
\usepackage{pdflscape}
\usepackage{graphicx}
\usepackage{tikz}
\usepackage{xcolor}
\usepackage[a4paper,top=2cm,bottom=2.5cm,left=2.5cm,right=2cm]{geometry}
\usepackage{fancyhdr}
\usepackage{draftwatermark}
\SetWatermarkText{La RevueCapes}
\SetWatermarkLightness{0.9}
\SetWatermarkScale{2}
\usepackage{hyperref}
\hypersetup{colorlinks=true,linkcolor=blue!50!black,urlcolor=blue!50!black}

\definecolor{revue}{RGB}{20,60,120}

\pagestyle{fancy}
\fancyhf{}
\renewcommand{\headrulewidth}{0.4pt}
\setlength{\headheight}{15pt}

% \entete{Matière}{Titre}{Type de document}
\newcommand{\entete}[3]{%
  \fancyhead[L]{\small\textsc{La Revue CAPES} -- #1}%
  \fancyhead[R]{\small #2}%
  \fancyfoot[C]{\thepage}%
  \thispagestyle{plain}%
  \begin{center}
    {\color{revue}\rule{\textwidth}{1.2pt}}\\[4pt]
    {\small\textsc{Concours directs CAP-CEG / CAPES -- Burkina Faso}}\\[6pt]
    {\LARGE\bfseries\color{revue} #2}\\[6pt]
    {\large\itshape #1 -- #3}\\[2pt]
    {\color{revue}\rule{\textwidth}{1.2pt}}
  \end{center}
  \vspace{0.5cm}%
}

% Encadré pour les remarques ajoutées dans les corrigés
\newcommand{\remarque}[1]{\par\medskip\noindent\fcolorbox{revue}{revue!6}{\parbox{\dimexpr\linewidth-2\fboxsep-2\fboxrule}{\textbf{\color{revue}Remarque.} #1}}\par\medskip}


\begin{document}
\entete{Mathématiques}{Corrigé détaillé du sujet type n°9}{Correction}
\begin{spacing}{1.5}

\textbf{Exercice 1}

\underline{Description de $D$.} Les courbes $x=\frac{y^2}{4}$ (parabole) et $y=2x$ (droite) se coupent lorsque $x=\frac{(2x)^2}{4}=x^2$, soit $x=0$ ou $x=1$ : aux points $O(0,0)$ et $(1,2)$. Pour $y\in[0,2]$, on a $\frac{y^2}{4}\le\frac y2$ (car $y(2-y)\ge0$). Donc
\[
D=\Big\{(x,y)\ /\ 0\le y\le2,\ \frac{y^2}{4}\le x\le\frac y2\Big\}=\Big\{(x,y)\ /\ 0\le x\le1,\ 2x\le y\le2\sqrt x\Big\}.
\]

\textbf{1. Aire de $D$.}
\[
\mathcal{A}(D)=\int_0^2\Big(\int_{y^2/4}^{y/2}dx\Big)dy=\int_0^2\Big(\frac y2-\frac{y^2}{4}\Big)dy=\Big[\frac{y^2}{4}-\frac{y^3}{12}\Big]_0^2=1-\frac{8}{12}=\boxed{\frac13}.
\]

\textbf{2. $\displaystyle\iint_D(y-x)\,dx\,dy$.} On intègre d'abord en $y$ :
\[
\int_{2x}^{2\sqrt x}(y-x)\,dy=\Big[\frac{y^2}{2}-xy\Big]_{2x}^{2\sqrt x}=(2x-2x\sqrt x)-(2x^2-2x^2)=2x-2x^{3/2},
\]
\[
\iint_D(y-x)\,dx\,dy=\int_0^1\big(2x-2x^{3/2}\big)dx=\Big[x^2-\frac45x^{5/2}\Big]_0^1=\boxed{\frac15}.
\]

\textbf{3. $\displaystyle\oint_Cy^2\,dx+xy\,dy$, $C$ orienté dans le sens direct (D à gauche).}

\textbf{a) Calcul direct.} Le bord $C$ est formé de deux arcs : le segment de droite (bord inférieur de $D$, parcouru de $O$ vers $(1,2)$) puis l'arc de parabole (bord supérieur, parcouru de $(1,2)$ vers $O$). On paramètre :
\[
\gamma_1:\ t\mapsto(t,2t),\quad t\in[0,1]\qquad\text{et}\qquad \gamma_2:\ t\mapsto(t,2\sqrt t),\quad t\in[0,1],
\]
$\gamma_2$ étant parcouru en sens inverse (de $t=1$ à $t=0$). Ainsi $\oint_C=\int_{\gamma_1}-\int_{\gamma_2}$.
\begin{itemize}
\item Sur $\gamma_1$ : $x=t$, $y=2t$, $dx=dt$, $dy=2dt$ ; $y^2dx+xy\,dy=4t^2dt+2t^2\cdot2\,dt=8t^2dt$, d'où $\displaystyle\int_{\gamma_1}=\int_0^18t^2dt=\frac83$.
\item Sur $\gamma_2$ : $x=t$, $y=2\sqrt t$, $dx=dt$, $dy=\frac{dt}{\sqrt t}$ ; $y^2dx+xy\,dy=4t\,dt+t\cdot2\sqrt t\cdot\frac{dt}{\sqrt t}=6t\,dt$, d'où $\displaystyle\int_{\gamma_2}=\int_0^16t\,dt=3$.
\end{itemize}
\[
\oint_Cy^2\,dx+xy\,dy=\frac83-3=\boxed{-\frac13}.
\]

\textbf{b) Par la formule de Green-Riemann.} Pour $P(x,y)=y^2$ et $Q(x,y)=xy$ (de classe $C^1$) :
\[
\oint_CP\,dx+Q\,dy=\iint_D\Big(\frac{\partial Q}{\partial x}-\frac{\partial P}{\partial y}\Big)dx\,dy=\iint_D(y-2y)\,dx\,dy=-\iint_Dy\,dx\,dy .
\]
Or
\[
\iint_Dy\,dx\,dy=\int_0^2y\Big(\frac y2-\frac{y^2}{4}\Big)dy=\Big[\frac{y^3}{6}-\frac{y^4}{16}\Big]_0^2=\frac86-1=\frac13 .
\]
On retrouve $\oint_Cy^2\,dx+xy\,dy=-\dfrac13$.

\medskip\textbf{Exercice 2}

\emph{Méthode générale (équations linéaires d'ordre 2 à coefficients constants).} La solution générale est $y=y_0+y_p$, où $y_0$ est la solution générale de l'équation homogène (déterminée par les racines de l'équation caractéristique) et $y_p$ une solution particulière. Pour un second membre $P(x)e^{mx}$ ($P$ polynôme), on cherche $y_p=x^kQ(x)e^{mx}$ avec $\deg Q=\deg P$ et $k\in\{0,1,2\}$ l'ordre de multiplicité de $m$ comme racine de l'équation caractéristique.

\textbf{1. $y''+2y'+y=2e^{-x}$, $y(0)=3$, $y'(0)=1$.}

\underline{Équation homogène.} $r^2+2r+1=(r+1)^2=0$ : racine double $r_0=-1$. D'où $y_0(x)=(Ax+B)e^{-x}$, $A,B\in\mathbb{R}$.

\underline{Solution particulière.} Le second membre est $2e^{-x}$ et $m=-1$ est racine \emph{double} : on cherche $y_p=Kx^2e^{-x}$.
\[
y_p'=K(2x-x^2)e^{-x},\qquad y_p''=K(2-4x+x^2)e^{-x},
\]
\[
y_p''+2y_p'+y_p=K\big[(2-4x+x^2)+2(2x-x^2)+x^2\big]e^{-x}=2Ke^{-x}.
\]
Donc $2K=2$, $K=1$ et $y_p=x^2e^{-x}$.

\underline{Solution générale et conditions initiales.} $y(x)=(x^2+Ax+B)e^{-x}$, et $y'(x)=\big(-x^2+(2-A)x+A-B\big)e^{-x}$.
$y(0)=B=3$ et $y'(0)=A-B=1$, d'où $A=4$ :
\[
\boxed{y(x)=(x^2+4x+3)e^{-x}=(x+1)(x+3)e^{-x}}.
\]

\textbf{2. Résolution des équations.}

\textbf{$(E_1)$ : $y''+y'=x+\operatorname{ch}x=x+\frac12e^x+\frac12e^{-x}$.}

\underline{Homogène.} $r^2+r=r(r+1)=0$ : $r_1=0$, $r_2=-1$. $y_0(x)=A+Be^{-x}$.

\underline{Particulières} (principe de superposition : on traite chaque terme du second membre).
\begin{itemize}
\item Second membre $x$ (polynôme, $m=0$ racine simple) : $y_1=x(ax+b)=ax^2+bx$ ; $y_1''+y_1'=2a+2ax+b=x$ donne $a=\frac12$, $b=-1$ : $y_1=\frac{x^2}{2}-x$.
\item Second membre $\frac12e^x$ ($m=1$ non racine) : $y_2=ce^x$ ; $y_2''+y_2'=2ce^x=\frac12e^x$ donne $c=\frac14$ : $y_2=\frac14e^x$.
\item Second membre $\frac12e^{-x}$ ($m=-1$ racine simple) : $y_3=dxe^{-x}$ ; $y_3'=d(1-x)e^{-x}$, $y_3''=d(x-2)e^{-x}$, $y_3''+y_3'=-de^{-x}=\frac12e^{-x}$ donne $d=-\frac12$ : $y_3=-\frac12xe^{-x}$.
\end{itemize}
\[
\boxed{y(x)=A+\Big(B-\frac x2\Big)e^{-x}+\frac{x^2}{2}-x+\frac14e^x,\qquad A,B\in\mathbb{R}}.
\]

\textbf{$(E_2)$ : $y''-3y'-4y=xe^{|x|}$.}

\underline{Homogène.} $r^2-3r-4=(r-4)(r+1)=0$ : $r_1=-1$, $r_2=4$. $y_0(x)=Ae^{-x}+Be^{4x}$.

Le second membre vaut $xe^{x}$ si $x\ge0$ et $xe^{-x}$ si $x<0$ ; il est continu sur $\mathbb{R}$. On résout sur chacun des intervalles $]0,+\infty[$ et $]-\infty,0[$, puis on raccorde en $0$.

\underline{Sur $]0,+\infty[$} : second membre $xe^x$, $m=1$ n'est pas racine. On cherche $y_+=(ax+b)e^x$ : $y_+'=(ax+a+b)e^x$, $y_+''=(ax+2a+b)e^x$ et
\[
y_+''-3y_+'-4y_+=\big[-6ax+(-a-6b)\big]e^x=xe^x\iff a=-\frac16,\ b=\frac1{36}.
\]
\underline{Sur $]-\infty,0[$} : second membre $xe^{-x}$, $m=-1$ est racine simple. On cherche $y_-=x(ax+b)e^{-x}=(ax^2+bx)e^{-x}$. On a
\[
y_-'=\big(-ax^2+(2a-b)x+b\big)e^{-x},\qquad y_-''=\big(ax^2+(b-4a)x+2a-2b\big)e^{-x},
\]
et $y_-''-3y_-'-4y_-=\big(-10ax+2a-5b\big)e^{-x}=xe^{-x}\iff a=-\frac1{10},\ b=\frac{2a}{5}=-\frac1{25}$.

Les solutions sur chaque intervalle sont donc
\[
y(x)=\begin{cases}
A_1e^{-x}+B_1e^{4x}+\Big(\dfrac1{36}-\dfrac x6\Big)e^x & \text{si } x>0,\\[3mm]
A_2e^{-x}+B_2e^{4x}-\Big(\dfrac{x^2}{10}+\dfrac{x}{25}\Big)e^{-x} & \text{si } x<0.
\end{cases}
\]
\underline{Raccordement en $0$.} Une solution sur $\mathbb{R}$ doit être deux fois dérivable en $0$. Il suffit que $y$ et $y'$ soient continues en $0$ : alors $y''=3y'+4y+xe^{|x|}$ a des limites égales à droite et à gauche, et $y$ est de classe $C^2$. Les valeurs en $0^+$ et $0^-$ sont :
\begin{align*}
y(0^+)&=A_1+B_1+\tfrac1{36}, & y(0^-)&=A_2+B_2,\\
y'(0^+)&=-A_1+4B_1-\tfrac5{36}, & y'(0^-)&=-A_2+4B_2-\tfrac1{25}.
\end{align*}
En égalant : $A_2-A_1+B_2-B_1=\frac1{36}$ et $-(A_2-A_1)+4(B_2-B_1)=-\frac5{36}+\frac1{25}$. En additionnant : $5(B_2-B_1)=-\frac4{36}+\frac1{25}=-\frac{16}{225}$, d'où
\[
B_2=B_1-\frac{16}{1125},\qquad A_2=A_1+\frac1{36}+\frac{16}{1125}=A_1+\frac{21}{500}.
\]
\textbf{Conclusion :} les solutions de $(E_2)$ sur $\mathbb{R}$ sont, pour $A,B\in\mathbb{R}$ :
\[
y(x)=\begin{cases}
Ae^{-x}+Be^{4x}+\Big(\dfrac1{36}-\dfrac x6\Big)e^x & \text{si } x\ge0,\\[3mm]
\Big(A+\dfrac{21}{500}\Big)e^{-x}+\Big(B-\dfrac{16}{1125}\Big)e^{4x}-\Big(\dfrac{x^2}{10}+\dfrac{x}{25}\Big)e^{-x} & \text{si } x<0.
\end{cases}
\]

\textbf{$(E_3)$ : $y''+y'-2y=x^2e^{-2x}$.}

\underline{Homogène.} $r^2+r-2=(r-1)(r+2)=0$ : $r_1=1$, $r_2=-2$. $y_0(x)=Ae^x+Be^{-2x}$.

\underline{Particulière.} $m=-2$ est racine simple. On pose $y=u(x)e^{-2x}$ : $y'=(u'-2u)e^{-2x}$, $y''=(u''-4u'+4u)e^{-2x}$, et
\[
y''+y'-2y=\big(u''-4u'+4u+u'-2u-2u\big)e^{-2x}=(u''-3u')e^{-2x}.
\]
Il suffit donc de trouver $u$ tel que $u''-3u'=x^2$. Comme il n'y a pas de terme en $u$, on cherche $u=ax^3+bx^2+cx$ :
\[
u''-3u'=6ax+2b-3(3ax^2+2bx+c)=-9ax^2+(6a-6b)x+(2b-3c)=x^2 .
\]
Par identification : $a=-\frac19$, $b=a=-\frac19$, $c=\frac{2b}{3}=-\frac2{27}$.
\[
\boxed{y(x)=Ae^x+\Big(B-\frac{x^3}{9}-\frac{x^2}{9}-\frac{2x}{27}\Big)e^{-2x},\qquad A,B\in\mathbb{R}}.
\]

\textbf{$(E_4)$ : $4xy''+2y'-y=0$.}

L'équation est linéaire du second ordre à coefficients variables ; le coefficient $4x$ de $y''$ s'annule en $0$. On résout sur $]0,+\infty[$ et sur $]-\infty,0[$, puis on raccorde.

\underline{Sur $]0,+\infty[$.} On pose $x=t^2$ ($t=\sqrt x>0$) et $y(x)=z(t)$. Comme $\frac{dt}{dx}=\frac1{2t}$ :
\[
y'(x)=\frac{z'(t)}{2t},\qquad y''(x)=\frac{1}{2t}\,\frac{d}{dt}\Big(\frac{z'(t)}{2t}\Big)=\frac{z''(t)}{4t^2}-\frac{z'(t)}{4t^3}.
\]
Alors $4xy''+2y'-y=4t^2\Big(\dfrac{z''}{4t^2}-\dfrac{z'}{4t^3}\Big)+\dfrac{z'}{t}-z=z''-z$. L'équation devient $z''-z=0$, d'où $z(t)=C_1\operatorname{ch}t+C_2\operatorname{sh}t$ et
\[
y(x)=C_1\operatorname{ch}\sqrt x+C_2\operatorname{sh}\sqrt x\qquad(x>0).
\]
\underline{Sur $]-\infty,0[$.} On pose $x=-t^2$ ($t=\sqrt{-x}>0$), $y(x)=z(t)$ ; le même calcul (avec $\frac{dt}{dx}=-\frac1{2t}$) donne $z''+z=0$, d'où
\[
y(x)=C_3\cos\sqrt{-x}+C_4\sin\sqrt{-x}\qquad(x<0).
\]
\underline{Raccordement en $0$.} Une solution sur $\mathbb{R}$ doit être deux fois dérivable en $0$.
\begin{itemize}
\item Continuité : $y(0^+)=C_1$ et $y(0^-)=C_3$, donc $C_1=C_3=y(0)$.
\item Dérivabilité à droite : $\dfrac{y(x)-C_1}{x}=C_1\dfrac{\operatorname{ch}\sqrt x-1}{x}+C_2\dfrac{\operatorname{sh}\sqrt x}{x}$. Le premier terme tend vers $\frac{C_1}{2}$ (car $\operatorname{ch}u-1\sim\frac{u^2}{2}$), le second est équivalent à $\frac{C_2}{\sqrt x}$ qui tend vers l'infini si $C_2\neq0$. Il faut donc $C_2=0$, et alors $y'_d(0)=\frac{C_1}{2}$.
\item De même à gauche : il faut $C_4=0$, et alors $y'_g(0)=\frac{C_1}{2}$.
\end{itemize}
Avec $C_2=C_4=0$ et $C_1=C_3=C$, on remarque que, pour tout réel $x$,
\[
y(x)=C\sum_{n=0}^{+\infty}\frac{x^n}{(2n)!}
\]
(car $\operatorname{ch}\sqrt x=\sum\frac{(\sqrt x)^{2n}}{(2n)!}$ pour $x\ge0$ et $\cos\sqrt{-x}=\sum\frac{(-1)^n(-x)^n}{(2n)!}=\sum\frac{x^n}{(2n)!}$ pour $x\le0$). C'est la somme d'une série entière de rayon infini : $y$ est de classe $C^\infty$ sur $\mathbb{R}$, et elle vérifie l'équation sur $\mathbb{R}^*$, donc aussi en $0$ par continuité.
\[
\boxed{y(x)=\begin{cases}C\operatorname{ch}\sqrt x & \text{si } x\ge0\\ C\cos\sqrt{-x} & \text{si } x<0\end{cases},\qquad C\in\mathbb{R}}.
\]
L'ensemble des solutions sur $\mathbb{R}$ est une droite vectorielle (et non un plan, à cause du point singulier $x=0$).

\medskip\textbf{Exercice 3}

\textbf{1. Pôle simple et coefficient $\dfrac{N(a)}{D'(a)}$.}

$F=\frac ND$ est irréductible : $N$ et $D$ n'ont pas de racine commune. Dire que $a$ est un pôle simple de $F$ signifie que $a$ est une racine simple de $D$ : $D(X)=(X-a)Q(X)$ avec $Q(a)\neq0$ (et $N(a)\neq0$).

\underline{$D'(a)\neq0$.} $D'(X)=Q(X)+(X-a)Q'(X)$, donc $D'(a)=Q(a)\neq0$.

\underline{Coefficient.} La décomposition en éléments simples de $F$ s'écrit
\[
F(X)=\frac{\lambda}{X-a}+G(X),
\]
où $G$ (partie entière et éléments simples relatifs aux autres pôles) n'a pas de pôle en $a$. En multipliant par $X-a$ :
\[
\frac{N(X)}{Q(X)}=\lambda+(X-a)G(X).
\]
En évaluant en $X=a$ (ce qui est licite car $Q(a)\neq0$ et $G$ est définie en $a$) : $\lambda=\dfrac{N(a)}{Q(a)}=\dfrac{N(a)}{D'(a)}$.

\emph{Exemple :} $F=\dfrac{1}{X^2-1}$, pôle $a=1$ : coefficient $\dfrac{1}{2\cdot1}=\dfrac12$.

\textbf{2. Négation.} La phrase s'écrit « Pour toute ville du Burkina, il existe au moins une agence bancaire dans cette ville ». Sa négation (on remplace « pour tout » par « il existe » et inversement, et on nie la propriété) est :
\begin{center}
\emph{« Il existe au moins une ville du Burkina qui ne possède aucune agence bancaire. »}
\end{center}

\textbf{3. Contraposée.} La phrase est de la forme $P\Rightarrow Q$ avec $P$ : « la pluviométrie n'est pas bonne » et $Q$ : « la campagne agricole ne sera pas une réussite ». La contraposée est $\operatorname{non}Q\Rightarrow\operatorname{non}P$ :
\begin{center}
\emph{« Si la campagne agricole est une réussite, alors la pluviométrie a été bonne. »}
\end{center}
\remarque{« Si la pluviométrie est bonne, alors la campagne sera une réussite » ($\operatorname{non}P\Rightarrow\operatorname{non}Q$) n'est pas la contraposée : c'est la contraire de la proposition, qui ne lui est pas équivalente.}

\textbf{4.a) $AB=0_n\Rightarrow A=0_n$ ou $B=0_n$ : FAUX.}

Contre-exemple pour $n=2$ : $A=\begin{pmatrix}1&0\\1&0\end{pmatrix}$ et $B=\begin{pmatrix}0&0\\1&1\end{pmatrix}$. Alors
\[
AB=\begin{pmatrix}1\cdot0+0\cdot1&1\cdot0+0\cdot1\\1\cdot0+0\cdot1&1\cdot0+0\cdot1\end{pmatrix}=\begin{pmatrix}0&0\\0&0\end{pmatrix},
\]
alors que $A\neq0_2$ et $B\neq0_2$. (L'anneau $\mathcal{M}_n(\mathbb{R})$, $n\ge2$, n'est pas intègre.)

\textbf{4.b) FAUX.} D'après les lois de De Morgan, $\operatorname{non}(P\text{ ou }Q)\iff(\operatorname{non}P)\text{ et }(\operatorname{non}Q)$. La négation de « $a\le b$ ou $|a|>c$ » est donc
\[
\text{« } a>b \textbf{ et } |a|\le c \text{ »},
\]
et non « $a>b$ ou $|a|\le c$ ».

\textbf{5. $f\big(A\cap f^{-1}(B)\big)=f(A)\cap B$.} On rappelle : $f^{-1}(B)=\{x\in E\ /\ f(x)\in B\}$.

\underline{Inclusion $\subset$.} Soit $y\in f\big(A\cap f^{-1}(B)\big)$. Il existe $x\in A\cap f^{-1}(B)$ tel que $y=f(x)$. Comme $x\in A$, $y=f(x)\in f(A)$ ; comme $x\in f^{-1}(B)$, $y=f(x)\in B$. Donc $y\in f(A)\cap B$.

\underline{Inclusion $\supset$.} Soit $y\in f(A)\cap B$. Comme $y\in f(A)$, il existe $x\in A$ tel que $y=f(x)$. Comme $f(x)=y\in B$, on a $x\in f^{-1}(B)$. Donc $x\in A\cap f^{-1}(B)$ et $y=f(x)\in f\big(A\cap f^{-1}(B)\big)$.

Par double inclusion, $f\big(A\cap f^{-1}(B)\big)=f(A)\cap B$.

\medskip\textbf{Exercice 4}

$f(z)=\dfrac{z-1}{z-2i}$ pour $z\neq2i$ ; $A(1)$, $B(2i)$, $M(z)$.

\textbf{1) Forme algébrique et forme trigonométrique de $f(i)$.}
\[
f(i)=\frac{i-1}{i-2i}=\frac{-1+i}{-i}=\frac{(-1+i)\,i}{-i\cdot i}=\frac{-i+i^2}{1}=\boxed{-1-i}.
\]
$|f(i)|=\sqrt{1+1}=\sqrt2$, et $f(i)=\sqrt2\Big(-\frac{\sqrt2}{2}-i\frac{\sqrt2}{2}\Big)$, donc
\[
\boxed{f(i)=\sqrt2\Big(\cos\Big(-\frac{3\pi}{4}\Big)+i\sin\Big(-\frac{3\pi}{4}\Big)\Big)=\sqrt2\,e^{-3i\pi/4}}.
\]

\textbf{2) Équation $f(z)=2i$.} Pour $z\neq2i$ :
\begin{align*}
\frac{z-1}{z-2i}=2i&\iff z-1=2i(z-2i)=2iz+4\iff z(1-2i)=5\\
&\iff z=\frac{5}{1-2i}=\frac{5(1+2i)}{5}=1+2i .
\end{align*}
Comme $1+2i\neq2i$, $\boxed{S=\{1+2i\}}$.

\textbf{3) Ensemble $D$ des points $M$ tels que $|f(z)|=2$.}

Pour $z\neq2i$ : $|f(z)|=2\iff|z-1|=2|z-2i|\iff AM=2BM$. Avec $z=x+iy$ :
\begin{align*}
(x-1)^2+y^2=4\big(x^2+(y-2)^2\big)&\iff3x^2+3y^2+2x-16y+15=0\\
&\iff x^2+y^2+\frac23x-\frac{16}{3}y+5=0 .
\end{align*}
En complétant les carrés :
\[
\Big(x+\frac13\Big)^2+\Big(y-\frac83\Big)^2=\frac19+\frac{64}9-5=\frac{20}{9}.
\]
Le point $B(0,2)$ n'est pas sur ce cercle ($\frac19+\frac49=\frac59\neq\frac{20}9$). Donc
\[
\boxed{D \text{ est le cercle de centre } \Omega\Big(-\frac13\,;\,\frac83\Big) \text{ et de rayon } \frac{2\sqrt5}{3}}.
\]
(C'est un cercle d'Apollonius : ensemble des points dont le rapport des distances à $A$ et à $B$ vaut 2.)

\end{spacing}
\end{document}
